\section{人类上下文窗口：人做熵减，AI 做熵增}

本章从开场的“长文本对话”开始。张小珺问李想，如果你是一个大模型，上下文窗口有多大。李想的回答没有沿着 token 数字走，而是转向人和 AI 的差别：人类不擅长处理特别复杂的信息，所以人要做熵减；AI 擅长吞吐巨量信息，更像做熵增。这里的“熵减”不是物理学严格定义，而是一种管理和认知方法：把复杂世界压缩成少数关键判断、少数动作和少数原则。

\begin{figure}[H]
\centering
\includegraphics[width=0.94\textwidth]{human-ai-context-window.png}
\caption{人类上下文 vs AI 上下文：人更适合做减法，AI 更适合处理大规模信息。自制概念图，依据 00:02:35--00:04:10 对谈内容整理。}
\end{figure}

\begin{knowledgebox}{读图：上下文窗口不是越大越像 CEO}
大模型可以吞吐很多 token，但 CEO 的价值常常在于压缩：把复杂信息变成战略、组织和行动。反过来，AI 如果只停留在“知道很多”，没有连接工具、行动和反馈，也不能产生真实价值。
\end{knowledgebox}

\subsection{工具三级：信息工具、辅助工具、生产工具}

上一节把人类与 AI 的上下文能力区分开，本节进一步追问：这种能力差异怎样变成产品价值。李想把工具分成三个层级。信息工具让人查到信息；辅助工具提高人的效率；生产工具则直接产生可付费的生产力。这个分类很重要，因为它给 AI 产品估值提供了一个更硬的标准：不是能不能回答，而是用户是否愿意为它付钱，是否愿意把真实工作交给它。

\begin{importantbox}{生产工具的判别标准}
如果一个 AI 工具只让人“觉得聪明”，它仍然可能只是信息工具；如果它能替代或解放人的高频真实工作，并且用户愿意为结果付钱，它才开始进入生产工具层级。
\end{importantbox}

\subsection{DeepSeek：极简地运用最佳实践}

接下来，访谈转向 DeepSeek。李想从 DeepSeek 学到的不是某个单点技巧，而是“极简地运用人类最佳实践”：先做研究，再做研发，再转化为业务；如果是业绩差距，就调整策略；如果是能力差距，就回到能力建设。这种路线看起来朴素，但很反人性，因为它要求组织不断承认自己能力不足，而不是只靠业务动作硬推。

\begin{figure}[H]
\centering
\includegraphics[width=0.96\textwidth]{deepseek-best-practice.png}
\caption{从 DeepSeek 学到的最佳实践：极简地使用人类最佳实践，再把组织和工程压实。自制概念图，依据 00:15:01--00:19:10 对谈内容整理。}
\end{figure}

\begin{knowledgebox}{读图：DeepSeek 被当作组织案例而不只是模型案例}
图里从 MoE 到实践的顺序，表达的是一套组织学习路径：先研究问题，再构建能力，再工程验证，最后让能力变成业务结果。李想关心 DeepSeek，是因为它展示了小团队、高密度人才、开源和工程实践如何形成合力。
\end{knowledgebox}

\begin{warningbox}{“最佳实践”并不自动产生最佳结果}
最佳实践往往反人性，因为它要求慢下来做研究、做能力、做复盘。随心所欲更满足人性，但容易把问题误判成短期策略问题，而忽略真正的能力缺口。
\end{warningbox}

\subsection{为什么还要自研基座能力}

访谈里有一个细节：理想内部也担心“用开源模型会不会伤害自研团队”。李想的回答是，开源语言模型可以成为一个基础，但理想要做的是基于业务场景发展出自己的 VLA、Agent OS 和物理世界能力。这里的核心不是“开源还是自研”的二选一，而是开源底座如何被嵌入公司自己的数据、任务、工具链和场景里。

\begin{knowledgebox}{开源模型与公司能力的关系}
Linux 可以开源，但安卓仍然形成了自己的组件、API、生态和产品体验。同理，语言模型可以开源，车企仍然需要围绕车、机器人、数据和用户场景构建自己的系统能力。
\end{knowledgebox}

\subsection{本章小结}

开场章节给整期定了基调：AI 不只是大上下文和大参数，而是怎样把信息处理能力变成行动能力。DeepSeek 被讨论，是因为它把研究、工程和组织压缩成了一条可学习的方法论。

